\section{1 绪论}

\subsection{1.1 研究背景与意义}

城市污水管网系统是城市基础设施的关键环节，承担着收集和输送生活污水、工业废水及雨水的重要功能。近年来，随着城镇化进程加快和污水收集范围不断扩大，管网系统中碳污染物的赋存与迁移问题逐渐增加，已成为制约城市碳减排和碳中和目标实现的关键因素之一。研究表明，污水管网同时承担碳污染物输送和生物转化的双重功能，一个复杂的生物化学反应器，管道内的微生物活动可导致碳污染物在固、液、气三相之间发生显著的相间迁移（Guisasola et al., 2008; Jiang et al., 2011）。

在"双碳"目标背景下，准确掌握污水管网中碳污染物的分布特征与迁移规律，对于城市碳排放核算、污水厂进水碳源优化以及管网运行管理对碳减排和管网管理具有实际价值。然而，现有研究多关注城市级污水管网系统，针对校园这一特殊功能区域的污水管网碳污染特征研究相对不足。校园污水管网具有排放源类型多样（教学区、生活区、餐饮区等）、用水规律性强、管道规模适中等特点，是研究碳污染物多相态分布特征的理想微尺度模型系统。

\subsection{1.2 国内外研究现状}

\subsubsection{1.2.1 污水管网碳污染物研究进展}

污水管网中碳污染物的研究始于20世纪90年代。早期研究主要关注管道中有机碳的生物转化过程，Guisasola等（2008）通过实验和模型模拟证实，污水在管道输送过程中，厌氧条件下的产甲烷活动可降解50\%以上的有机碳，使管道成为一个重要的碳转化器。Jiang等（2011）进一步指出，管道系统中产生的CH$\_4$和CO$\_2$是城市温室气体排放的重要来源，其排放量占城市碳排放总量的显著比例。

近年来，多相态分析方法逐渐被引入污水管网碳污染物研究。研究者开始同时关注固相（管道沉积物和生物膜中的有机碳）、液相（溶解性有机碳DOC和颗粒态有机碳POC）和气相（CH$\_4$和CO$\_2$）碳污染物的分布特征及其相互转化关系。然而，多数研究仅关注单一相态，缺乏对固-液-气三相碳污染物的系统性联合分析。

\subsubsection{1.2.2 校园污水管网研究现状}

校园污水管网研究尚处于起步阶段。现有少量研究主要集中在水质指标的基础监测层面，对碳污染物在管网中的相态分布、空间分异及其影响因素缺乏分析。校园污水管网因其排放源明确、空间尺度适中、易于系统采样等优势，可作为研究污水管网碳转化过程的理想实验平台。

\subsection{1.3 现有研究不足}

，现有研究存在以下不足：

（1）\textbf{多相态联合分析不足。} 已有研究多关注单一相态的碳污染物，缺乏固-液-气三相碳污染物的系统性联合分析，难以全面揭示碳在管网中的分布特征和相间迁移规律。

（2）\textbf{校园尺度研究缺乏。} 现有研究以城市级管网为主，针对校园这一特殊功能区域的碳污染特征研究较少，对不同功能区（教学区、生活区、餐饮区）碳排放差异的认识不足。

（3）\textbf{碳平衡分析薄弱。} 管网中碳的输入-输出平衡关系尚不清楚，碳在不同相态之间的分配比例及其影响因素有待定量揭示。

（4）\textbf{影响因素不明。} 溶解氧、温度、有机负荷等因素对碳污染物相间迁移的影响因素缺乏研究。

\subsection{1.4 研究内容与目标}

针对上述研究不足，本研究以某校园污水管网为研究对象，开展以下研究工作：

（1）系统采集管道内固相、液相、气相样品，测定碳污染物各指标浓度，揭示固-液-气三相碳污染物的分布特征。

（2）采用主成分分析(PCA)和层次聚类分析(HCA)等多元统计方法，识别影响碳污染物分布的关键因素和采样点聚类特征。

（3）分析不同功能区碳污染物的空间分异规律，探讨排放源类型对碳污染特征的影响。

（4）开展碳平衡分析，定量评估碳在固-液-气三相之间的分配比例及其影响因素，为校园污水碳管理提供数据支撑。